\chapter{Conclusiones}
\label{ch:conclusiones}

La resolución del Ejercicio Práctico N° 4 permitió al grupo comprender e implementar los conceptos fundamentales
de la comunicación serie asíncrona bidireccional mediante el periférico USART1 en microcontroladores de arquitectura
Cortex-M3.

A lo largo del desarrollo se presentaron desafíos técnicos significativos que enriquecieron el aprendizaje del grupo.
El primer obstáculo importante radicó en la configuración del divisor de baudios (registro \lstinline{USART1->BRR}).
Inicialmente, se asumió que la frecuencia de reloj del sistema era de \qty{72}{\MHz} de acuerdo a la configuración
máxima típica del microcontrolador. Al aplicar la máscara correspondiente a dicha frecuencia, la comunicación
fracasó de manera sistemática y la PC sólo recibía bytes nulos (\texttt{0x00}) o ruidos de sincronismo. La depuración
de este problema forzó al grupo a examinar minuciosamente el datasheet, confirmando que el microcontrolador se mantiene funcionando con el oscilador interno
HSI de \qty{8}{\MHz} en su estado por defecto, requiriendo recalcular el registro divisor a \lstinline{0x341} para
lograr la sincronización a $9600$ baudios.

Otro error frecuente y silencioso ocurrió al intentar integrar el temporizador SysTick para la gestión de retardos y el
parpadeo del LED. La función de interrupción se nombró originalmente como \lstinline{SysTick_IRQHandler}. Dado que
el enlazador de GCC no encontró errores de sintaxis, el código compilaba correctamente. Sin embargo, el contador
\lstinline{ticks} no se incrementaba en absoluto. Tras consultar la tabla de vectores en el código de startup en
ensamblador (\texttt{startup\_stm32f103xb.s}), se descubrió que el nombre exacto de la función en la tabla es
\lstinline{SysTick_Handler} (sin la sigla IRQ). Al corregir este identificador, el enlazador sustituyó exitosamente la
referencia débil y la base de tiempo de \qty{1}{\ms} comenzó a funcionar de forma transparente.
